% Department cover of the capstone report (official format: A4, 12 pt, margins 2.54 / 3.17 cm). build.sh compiles it
% and puts it in front of the IEEE two-column body (report.tex), as the 2025 winners did.
\documentclass[12pt,a4paper]{article}
\usepackage[top=2.54cm,bottom=2.54cm,left=3.17cm,right=3.17cm]{geometry}
\usepackage{fontspec}
\usepackage{xeCJK}
\setmainfont{Liberation Serif}
\setCJKmainfont{Noto Serif TC}
\usepackage{graphicx}
\usepackage{booktabs}
\usepackage{adjustbox}
\usepackage{array}
\usepackage{tabularx}
\usepackage{xcolor}
\usepackage{tikz}
\usetikzlibrary{positioning,arrows.meta,calc,shapes.geometric}
\usepackage{pifont}
\usepackage{caption}
\usepackage{enumitem}
\usepackage{amsmath}
\newcommand{\member}[1]{\textcolor{blue}{[#1]}}
\pagestyle{empty}
\begin{document}
\begin{titlepage}
\centering
{\large 國立清華大學資訊工程系 115 學年度上學期專題報告}\\[0.3em]
{\normalsize National Tsing Hua University, Department of Computer Science, Undergraduate Project Report, Fall 2026}\\[1.6em]
\renewcommand{\arraystretch}{1.5}
\begin{tabularx}{\textwidth}{|>{\centering\arraybackslash}p{3.2cm}|X|}
\hline
專題名稱 & \textbf{Beautified Is Not Forged: A Face Checker That Separates Beautification from Forgery and Shows Checkable Evidence} \\
\hline
參加競賽或計畫 & \member{□ 參加對外競賽　□ 參與其他計畫　□ 無參加對外競賽或任何計畫} \\
\hline
\end{tabularx}\\[0.8em]
\begin{tabularx}{\textwidth}{|>{\centering\arraybackslash}p{3.2cm}|X|X|X|}
\hline
學號 & \member{student ID} & \member{student ID} & \member{student ID} \\
\hline
姓名 & 林沁瑩 & 陳柏蓉 & 余晨杉 \\
\hline
\end{tabularx}\\[0.8em]
\begin{tabularx}{\textwidth}{|>{\centering\arraybackslash}p{3.2cm}|X|}
\hline
摘要 &
\small % Abstract, shared by cover.tex and report.tex.
Face-forgery detectors answer ``real'' or ``fake''. A genuine photograph that has only been through a beauty app or a
retouching service is neither, and once the skin is smoothed the detectors tend to call it fake. On pixel-aligned pairs of
genuine faces and their commercially smoothed renders, sixteen of seventeen published detectors score the smoothed face as more
forged than its own original. A controlled training experiment shows that the cause is which class saw skin smoothing during
training: showing it to both classes lowers the accusation from 54 to 12 percent. We build on that recipe a face checker that
answers real, forged or beautified, marks the face part its evidence came from, and shows only sentences that code composes from
the model's output and checks before display. Restoring the marked region to its original pixels removes 94 percent of the
model's part-level forgery decisions, about twice what a post-hoc heat map achieves, and free-form descriptions written by
vision-language models fail a blind matching test, so no model is allowed to describe the image. The checker runs in a web
browser and the photograph never leaves the device.
 \\
\hline
\end{tabularx}\\[2em]
{\large 中華民國 115 年 11 月}
\end{titlepage}
\end{document}
